\documentclass[11pt]{article}
\usepackage[margin=0.9in]{geometry}
\usepackage{amsmath,amssymb,amsthm,booktabs,mathrsfs}
\usepackage[T1]{fontenc}
\usepackage{lmodern,microtype}
\usepackage[colorlinks=true,linkcolor=blue,urlcolor=blue]{hyperref}
\newtheorem{theorem}{Theorem}[section]
\newtheorem{lemma}[theorem]{Lemma}
\newtheorem{proposition}[theorem]{Proposition}
\newtheorem{corollary}[theorem]{Corollary}
\theoremstyle{remark}
\newtheorem{remark}[theorem]{Remark}
\newcommand{\norm}[1]{\left\|#1\right\|}
\title{Toward a four-dimensional Yang--Mills construction\\
\large Notebook 19: the physical background Gaussian in the source coordinates}
\author{Research working document}
\date{17 October 2026}
\begin{document}
\maketitle
\begin{abstract}
We construct the actual homogeneous current-step pivot chart near
the trivial field and a physical stationary background for the
next ordinary averaging constraint. The chart is local and leaves
the retained coordinates unchanged. Two minimum comparisons use
the previously constructed ordinary minimizer to identify the
current classical action on the next fiber. An explicit rooted
gauge sets the source's averaged contour variable to the identity.
The resulting reduced fluctuation Hessian is a small perturbation
of the flat retained precision, uniformly in earlier averaging
depth and periodic volume. Its formula keeps the nonlinear
constraint curvature and the conversion from relative exponential
to ordinary potential coordinates. This gives a uniform complex
covariance neighborhood for the physical small-background branch.
We also isolate the local measure and gauge factors needed by
the interacting integral.
\end{abstract}

\section{From a classical minimum to the source Gaussian}
\label{ps19-target}

The existence argument needs the Gaussian in the actual quantum
integration coordinates. The fine constrained inverse and the
Hessian of a minimized coarse action are useful inputs, but
neither alone identifies that Gaussian. The source also changes
coordinates, eliminates pivot bonds, and includes an exponential
gauge term. Away from the flat field these operations contribute
current-dependent second derivatives.

Here the current lattice has unit spacing, the previous fine
spacing is $M^{-1}$, and the next spacing is a fixed odd $L$.
The background is in a sufficiently small global potential ball.
The constants are uniform in $M=L^k$ and compatible periodic
volume, with the specified source block/interface geometry.
The result concerns physical stationary backgrounds and their
analytic extension near the trivial field. The original larger
orbit domain, freely prescribed independent currents, and mixed
large-field histories require further estimates.

